\documentclass[12pt,a4j]{jreport}
\usepackage[dvipdfmx]{graphicx}
\begin{document}
\thispagestyle{empty}
\begin{center}
修士論文/ 課題研究報告書\\% どちらか一方を消してください．
\vfill
〇〇〇〇〇題目〇〇〇〇〇\\
\vfill
〇〇著者名〇〇\\
\vfill
主指導教員  〇〇 〇〇\\
\vfill
北陸先端科学技術大学院大学\\
先端科学技術研究科\\
（〇〇〇〇）\\ %取得希望学位
\vfill
令和〇〇年〇月\\ % 学位授与年月
\vfill
\end{center}
\clearpage
\pagestyle{empty}
\centerline{Abstract}
English abstract (1200 words)\\
〇〇〇〇〇〇〇〇〇〇〇〇〇〇〇〇〇〇〇
\clearpage
\centerline{概要}
日本語の概要\\
〇〇〇〇〇〇〇〇〇〇〇〇〇〇〇〇〇〇〇

\clearpage\setcounter{page}{0}\pagenumbering{roman}\pagestyle{plain}
\tableofcontents
\listoffigures
\listoftables
\clearpage\setcounter{page}{0}\pagenumbering{arabic}

\chapter{はじめに}
〇〇〇〇〇〇〇〇〇〇〇〇〇〇〇〇〇〇\cite{jref1,jref2,jref3}
\chapter{関連研究}
\section{セクション名}
\subsection{サブセクション名}
〇〇〇〇〇〇〇〇〇〇〇〇〇〇〇〇〇〇
\chapter{提案手法}
〇〇〇〇〇〇〇〇〇〇〇〇〇〇〇〇〇〇 (図 \ref{fig1})
\begin{figure}
\centering
% \includegraphics[clip,width=90mm]{boltzmann.jpg}
(図を貼る)
% 図のキャプションは図の下に置く。
\caption{図のキャプション}\label{fig1}
\end{figure}
\chapter{実験・評価}
\begin{table}
\centering
% 表のキャプションは表の上に置く。
\caption{表のキャプション}\label{table1}
\begin{tabular}{r|rr}
& a & b\\ \hline
1& 0.25 & 0.33\\
2& 0.75 & 0.66\\
\end{tabular}
\end{table}
\chapter{おわりに}

% 参考文献
\renewcommand{\bibname}{参考文献}

% BibTeX を使わないとき
% \renewcommand{\bibname}{参考文献}
% \begin{thebibliography}{99}
% \bibitem{jref1} 著者名. 本のタイトル. 出版社, 出版年.
% \bibitem{jref2} 著者1 and 著者2. 論文タイトル. 学会名, pp.zz--ww, 2007.
% \bibitem{jref3} 著者1, 著者2, and 著者3. 論文タイトル. ジャーナル名, vol.xx, no.yy, pp.zz--ww, 2003.
% \end{thebibliography}

% BibTeX を使うとき
\renewcommand{\em}{\it}% booktitle に下線を引かない
\bibliographystyle{jplain}
\bibliography{reference}

% 付録
% \appendix
% \chapter{参考データ}

\end{document}
